\subsubsection*{1.5.}
设 $\bar\eta$ 为 $S$ 的一个几何点 (位于 $S$ 的点 $\eta$ 之上), 并设
\begin{equation}\tag{1.5.1}
z_{\bar\eta}\in
H^a\bigl(X_{\bar\eta}\times_{k(\bar\eta)}Y_{\bar\eta},
\mathbb Q_\ell(b)\bigr)
\end{equation}
为 $z$ 在
\begin{equation}\tag{1.5.2}
X\times_S Y\longrightarrow S.
\end{equation}
的相应几何纤维上诱导的类.

\subsubsection*{命题 1.5.3.}
假定 $f$ 和 $g$ 固有, $S$ 连通,\footnote{印本以无意义的形容词 ``annexe'' 代替 ``connexe'', 并将结论中的元素写作大写 $X$. 证明要求 $S$ 连通, 而 (1.5.4) 本身使用元素 $x$.} 并设 $q\in\mathbb Z$. 若类 $z_{\bar\eta}$ 诱导零同态
\begin{equation}\tag{1.5.3}
0=z_{\bar\eta}:H^{q-1}(X_{\bar\eta},\mathbb Q_\ell(i))
\longrightarrow
H^{q+a-2d_1-1}(Y_{\bar\eta},\mathbb Q_\ell(i+b-d_1)),
\end{equation}
则对每个 $x\in\mathrm{Prim}^q(X/S,\mathbb Q_\ell(i))$, 有
\begin{equation}\tag{1.5.4}
u_g(z(x))=0\quad\text{于}\quad
H^1\bigl(S,R^{q+a-2d_1-1}g_*\mathbb Q_\ell(i+b-d_1)\bigr).
\end{equation}

\noindent\underline{证明.}
由 (1.4.2) 的交换性, 只须证明 $z$ 诱导零同态
\begin{equation}\tag{1.5.5}
z:R^{q-1}f_*\mathbb Q_\ell(i)
\stackrel{0}{\longrightarrow}
R^{q+a-2d_1-1}g_*\mathbb Q_\ell(i+b-d_1).
\end{equation}
事实上, 由于 $f$ 和 $g$ 固有且光滑, 层 $R^\bullet f_*\mathbb Q_\ell(\mathord\cdot)$ 和 $R^\bullet g_*\mathbb Q_\ell(\mathord\cdot)$ 为扭常层 (SGA~4, XVI~2.2), 因而等价于它们在 $\bar\eta$ 处的值, 其中将该值视为带有 $\pi_1(S,\bar\eta)$ 作用的模.

\subsubsection*{1.6. Manin 的一个构造}
假定给定 $f:X\longrightarrow S$ 的两个截面:
\begin{equation}\tag{1.6.0}
\begin{tikzcd}[column sep=huge]
S \arrow[r, bend left=45, "t_1"]
  \arrow[r, bend right=45, "t_2"']
& X \arrow[l, "f" description]
\end{tikzcd}
\qquad ,
\end{equation}
并赋予每个截面其上同调类
\begin{equation}\tag{1.6.1}
\mathrm{cl}(t_i)=t_{i*}(1)
\in H^{2d_1}(X,\mathbb Q_\ell(d_1)).
\end{equation}
若 $f:X\longrightarrow S$ 的纤维几何连通 (更一般地, 若对 $s\in S$, $t_1(s)$ 和 $t_2(s)$ 属于 $X_s$ 的同一个连通分支), 则显然有\footnote{印本此处写作 $f(s)$ 和 $g(s)$, 但正在讨论的两个映射是截面 $t_1$ 和 $t_2$.}
\begin{equation}\tag{1.6.2}
\mathrm{cl}(t_1)-\mathrm{cl}(t_2)
\in\mathrm{Prim}^{2d_1}(X/S,\mathbb Q_\ell(d_1)).
\end{equation}
以下式表示
\begin{equation}\tag{1.6.3}
u_f(t_1-t_2)\in
H^1(S,R^{2d_1-1}f_*\mathbb Q_\ell(d_1))
\end{equation}
即其经
\begin{equation}\tag{1.6.4}
u_f:\mathrm{Prim}^{2d_1}(X/S,\mathbb Q_\ell(d_1))
\longrightarrow H^1(S,R^{2d_1-1}f_*\mathbb Q_\ell(d_1)).
\end{equation}
所得的像.
(Manin 在~[5] 中研究的正是这个元素的一个类似物 $\mu(t_1-t_2)$.)

\subsubsection*{1.6.5.}
现在设 $z$ 为 $X\times_S Y$ 上的代数圈, 在 $X\times_S Y$ 的每个纤维上的余维至少为 $b$,\footnote{若 $S$ 在概形 $T$ 上光滑 (例如 $T=\operatorname{Spec}k$, 其中 $k$ 为域), 则只须假设 $z$ 在 $T$ 上的每个纤维中的余维至少为 $b$.} 它确定一个上同调类 (SGA~5, 第~IV 讲; 印本中后续的引用留空)
\begin{equation}\tag{1.6.6}
\mathrm{cl}(z)\in H^{2b}(X\times_S Y,\mathbb Q_\ell(b)).
\end{equation}
此类定义映射 (参见~(1.2.3))
\begin{equation}\tag{1.6.7}
\mathrm{cl}(z)\mathbin{\cdot}H^{2d_1}(X,\mathbb Q_\ell(d_1))
\longrightarrow H^{2b}(Y,\mathbb Q_\ell(b)),
\end{equation}
若圈 $z(t_1)$ 和 $z(t_2)$ 有定义, 则杯积与圈的相交的相容性 (SGA~5, 第~IV 讲; 印本中后续的引用留空) 给出
\begin{equation}\tag{1.6.8}
\mathrm{cl}(z)\bigl(\mathrm{cl}(t_1)-\mathrm{cl}(t_2)\bigr)
=\mathrm{cl}\bigl(z(t_1)-z(t_2)\bigr).
\end{equation}
由 (1.4.2) 的交换性得到
\begin{equation}\tag{1.6.9}
\mathrm{cl}(z)\bigl(u_f(t_1-t_2)\bigr)
=u_g\bigl(\mathrm{cl}(z(t_1)-z(t_2))\bigr).
\end{equation}
所以作为命题~1.5.3 的推论, 有:

\subsubsection*{命题 1.7.}
设 $t_1$ 和 $t_2$ 为 $f:X\longrightarrow S$ 的两个截面, 假定 $f$ 的纤维连通, 并设 $z$ 为 $X\times_S Y$ 上余维为 $b$ 的代数圈, 使得代数圈
\begin{equation}\tag{1.7.1}
z(t_1)-z(t_2)
\end{equation}
有定义. 则其上同调类为本原类 (XVIII~5.8.1). 若 $X_{\bar\eta}\times Y_{\bar\eta}$ 上的代数圈 $z_{\bar\eta}$ 诱导零同态
\begin{equation}\tag{1.7.2}
z_{\bar\eta}:H^{2d_1-1}(X_{\bar\eta},\mathbb Q_\ell(d_1))
\stackrel{0}{\longrightarrow}
H^{2b-1}(Y_{\bar\eta},\mathbb Q_\ell(b)),
\end{equation}
则
\begin{equation}\tag{1.7.3}
u_g\bigl(z(t_1)-z(t_2)\bigr)=0\quad\text{于}\quad
H^1(S,R^{2b-1}g_*\mathbb Q_\ell(b)),
\end{equation}
其中 $u_g$ 为~1.4 中出现的 XVIII~5.8.1 的同态.
\footnote{印本的句子中使用大写 $Z_{\bar\eta}$, 但 (1.7.2) 和本节始终使用的圈都记为 $z_{\bar\eta}$.}
